% =============================================================================
% SECTION 8
% =============================================================================
\newpage
\section{Conclusion}
\label{sec:conclusion}

\revised{\textbf{Scope.} The dual-jurisdiction analysis establishes functional equivalents under English, Delaware, and California law; conflicts and forum rules govern enforcement, while regulation and voluntary implementation provide parallel deployment routes. The framework is instrument-agnostic by design: psychometric tools calibrate governance thresholds to deployment context, but neither the selected instrument nor its numerical cut-off alters the doctrinal principle that protection scales with irreversible investment.}

\bigskip

On a Tuesday in early 2024, a company pushed a model update and something was taken from its users. This paper has argued that the something has a name. It was not necessarily a copyrighted work: every prompt and output might remain in the archive. It was the productive capacity of the personalised state to respond in the register those interactions had cultivated, and the user's periodic estate in that capacity. The update wasted the asset and displaced the tenant.

The framework proposed here requires no further legislation. The Property (Digital Assets etc) Act 2025 confirms that digital form is not a categorical bar to personal-property status; established estate doctrine supplies the rest. At Stage~1, tenancy, notice, and waste protect productive capacity without a trust. At Stage~2, deliberately cultivated dependence and unconscionable exploitation may justify the constructive trust and its fiduciary remedies.

The tenancy does the ordinary work. The trust is reserved for the extraordinary case the company itself cultivated.
